\uuid{Ua6q}
\exo7id{1620}
\titre{exo7 1620}
\auteur{liousse}
\organisation{exo7}
\datecreate{2003-10-01}
\isIndication{false}
\isCorrection{false}
\chapitre{Réduction d'endomorphisme, polynôme annulateur}
\sousChapitre{Diagonalisation}
\cpge{mpsi-mp2i.determinants.d, pcsi.calcul-matriciel-systemes-lineaires.d, ptsi.calcul-matriciel-systemes-lineaires.d, mpsi-mp2i.determinants.e, pcsi.determinants.d, mp-mpi.reduction-endomorphismes.b, pc.reduction.a, psi.reduction-endomorphismes.a, pt.reduction-endomorphismes.a, mp-mpi.reduction-endomorphismes.c, pc.reduction.b, psi.reduction-endomorphismes.b, pt.reduction-endomorphismes.b, mp-mpi.reduction-endomorphismes.d, pc.reduction.c, psi.reduction-endomorphismes.c, pt.reduction-endomorphismes.c, bcpst2.valeurs-propres-vecteurs-propres.a, bcpst2.valeurs-propres-vecteurs-propres.b}
\module{Algèbre}
\niveau{L2}
\difficulte{}

\contenu{
\question{
On consid\`ere la matrice 
$$M=\left (\begin{array}{ccc}
3&2&1\\-1&0&-1\\-1&-1&1\end{array}\right )$$ 
Est-elle diagonalisable ? Justifier. \'Ecrire alors $M$ sous 
une forme plus simple.
}
}
